\section{Introduction}\label{sec:introduction}

Laughlin's wave function describes a strongly correlated state at
fractional filling of the lowest Landau level~\cite{Laughlin}.
Haldane's spherical geometry and pseudopotentials give a particularly
simple parent Hamiltonian: each electron pair is penalized by the
orthogonal projection onto relative angular momentum one~\cite{Haldane}.
The cubic Laughlin polynomial is its unique zero state at filling
$1/3$ with the spherical shift. A spectral gap separates this state
from excitations. Stability asks whether that separation survives
a weak external potential whose total many-particle norm grows with
the area of the system.

We prove uniform stability for the full spherical interaction under
bounded scalar potentials projected to the lowest Landau level.
The potential may vary with system size and need not have any symmetry.
The proof uses the uniform unperturbed Fock-space gap of~\cite{Gap},
and derives the additional local estimates needed to pass from a
gap at zero disorder to a gap on a fixed disorder interval.

\subsection{Model and main result}

For $N\ge2$, set $q=3(N-1)$ and let $S_q$ be the round sphere of
radius $\sqrt{q/2}$. Magnetic length is one, so the sphere carries
$q$ flux quanta and has area $2\pi q$. Let $U_q$ be its one-particle
lowest Landau level, equivalently the spin-$q/2$ representation of
$SU(2)$. On the entire antisymmetric space $\bigwedge^N U_q$, put
\begin{equation}\label{eq:intro-H}
 H_{N,q}=\sum_{1\le i<j\le N}P_{ij}^{(1)},
\end{equation}
where $P_{ij}^{(1)}$ projects the indicated pair onto total spin
$q-1$. Each unordered pair has coefficient one. The normalized
Laughlin vector spans the line generated by
\begin{equation}\label{eq:intro-laughlin}
 \Psi_{\mathrm L,N}=\prod_{i<j}
       (z_{i0}z_{j1}-z_{i1}z_{j0})^3.
\end{equation}

If $\Pi_q$ denotes the lowest-level projection and $\varphi$ is a
bounded real scalar function on $S_q$, its Toeplitz operator and
many-particle perturbation are
\[
 T_q(\varphi)=(\Pi_qM_\varphi\Pi_q)|_{U_q},\qquad
 V_{N,q}(\varphi)=\sum_{i=1}^N T_q(\varphi)^{(i)}.
\]
Let $E_0(N,\lambda,\varphi)\le E_1(N,\lambda,\varphi)$ be the
two lowest eigenvalues, counted with multiplicity, of
\[
 H_{N,q}(\lambda,\varphi)=H_{N,q}+\lambda V_{N,q}(\varphi).
\]

\begin{theorem}[Uniform scalar-disorder stability]\label{thm:main}
There are constants $\lambda_*>0$, $\Delta_*>0$, and
$N_*<\infty$ such that, for every $N\ge N_*$, every real
bounded measurable $\varphi$ on $S_{3(N-1)}$ with
$\|\varphi\|_\infty\le1$, and every real
$|\lambda|\le\lambda_*$,
\begin{equation}\label{eq:intro-stability}
 E_1(N,\lambda,\varphi)-E_0(N,\lambda,\varphi)\ge\Delta_*.
\end{equation}
In particular the perturbed ground state is unique.
\end{theorem}

The constants are uniform over the potential as well as particle number.
Thus the theorem applies, in particular, to smooth potentials satisfying
$\max_{0\le k\le2}\|\nabla^k\varphi\|_\infty\le1$ in the
physical round metric. Its stronger amplitude-only formulation comes
from the coherent-orbital representation of Toeplitz multiplication.
The ground vector is allowed to move, and its energy is subtracted in
the stability argument. Adding a constant to the potential changes
only that energy. We give existential constants; the proof does not
optimize the disorder threshold.

\subsection{Historical context and the additional difficulty}

Exact parent Hamiltonians have been central to the mathematical study
of fractional quantum Hall states since the work of Trugman and
Kivelson~\cite{TK}. Haldane pseudopotentials also arise from suitable
short-range limits of repulsive interactions, as proved by Seiringer
and Yngvason~\cite{SY}. The spectral gap is a separate quantitative
question: identifying a zero polynomial does not control the energy
of vectors perpendicular to it. Rougerie's survey~\cite{Rougerie}
discusses the gap problem and its role in the response of the Laughlin
phase to perturbations.

Rigorous uniform gaps were established for truncated pseudopotentials
in thin-cylinder and thin-torus geometries by Nachtergaele, Warzel,
and Young~\cite{NWY} and Warzel and Young~\cite{WY}. The unperturbed
input used here concerns the full interaction on expanding round
spheres~\cite{Gap}. We state its exact normalization in
Theorem~\ref{thm:unperturbed}. In addition, we use and reproduce a
strengthening of its four-body comparison: certain positive blocks
that were discarded in obtaining the gap can be retained. This
stronger estimate supplies the energy-loss construction below.

Another line of work controls density within the Laughlin-correlated
class. Rougerie and Yngvason~\cite{RY}, followed by Lieb, Rougerie,
and Yngvason~\cite{LRY}, established incompressibility estimates for
Laughlin functions multiplied by analytic symmetric factors. The
sharp local density bound in~\cite{LRY} is averaged over mesoscopic
disks. Our zero-mode estimate controls general local observables,
with total cost proportional to the number of missing electrons;
it is then combined with an estimate for states outside the zero space.
Algebraic descriptions and recursions for pseudopotential zero modes
were developed in~\cite{CS,MONS}. We use the symmetric-polynomial
description and add quantitative bounds for its orthogonal branching.

General gap-stability theorems explain the role of local information.
Bravyi, Hastings, and Michalakis~\cite{BHM} treated stability under
sufficiently weak, suitably decaying local perturbations of
commuting-projector models with topological-order conditions. Michalakis and
Zwolak~\cite{MZ} proved stability for frustration-free lattice
Hamiltonians under local topological-order and local-gap hypotheses.
Subsequent infinite-volume formulations also retain local conditions
in addition to a bulk gap~\cite{NSYBulk}.
Those hypotheses contain more information than a uniform global gap.
For the spherical lowest Landau level, the local control needed here
is established through flux pinning and a particle-loss evolution.
The resulting relative-bound argument has the same purpose as local
block diagonalization in the lattice theory, with a different source
of control for the diagonal terms.

Our final continuation of the ground line uses the quasiadiabatic
construction of Hastings and Wen~\cite{HW} and the exact spectral-flow
framework of Bachmann, Michalakis, Nachtergaele, and Sims~\cite{BMNS}.
It is applied inside a gap bootstrap, where its gap hypothesis is
available. Locality is obtained using redundant localized frames.
This connects to the continuum fermionic frame method of Bachmann and
De Nittis~\cite{BDN}; we prove the finite-sphere estimates, including
comparisons at neighboring fluxes, that the present argument requires.
The propagation estimates follow the Lieb--Robinson method~\cite{LR}
and its fermionic form~\cite{NSY}.

\subsection{The proof and its reusable estimates}

The main difficulty is extensive perturbation size:
$\|V_{N,q}(\varphi)\|$ can be of order $N$. A norm bound on the
whole perturbation therefore gives a disorder interval that shrinks
with particle number. We instead compare suitably centered local
terms with the interaction energy. Two estimates make that comparison
possible.

\paragraph{Local observables in zero spaces.}
At fixed flux $q=3(N-1)$, write $Z_{n,q}=\ker H_{n,q}$.
For an interaction $W$ that is a bounded-density sum of uniformly
localized neutral terms, we prove
\begin{equation}\label{eq:intro-zero-bound}
 \left|\langle\psi,W\psi\rangle
       -\langle\Omega,W\Omega\rangle\right|
 \le C_W(N-n),\qquad
 \psi\in Z_{n,q},\quad\|\psi\|=1,
\end{equation}
where $\Omega$ is the normalized filled Laughlin vector.
This is Theorem~\ref{thm:charge-local-linear}; the constant is uniform
in flux. Neutral means that a term preserves
particle number; localization is defined precisely in
Section~\ref{sec:locality}. The linear dependence on $N-n$ is
essential.

To prove \eqref{eq:intro-zero-bound}, we read a zero mode from one
pole. Either the last orbital is empty, which removes one unit of
flux and records a flag, or contraction in that orbital removes an
electron and three units of flux. The branches are orthogonal and
give coordinates for the entire zero space. The number of flags is
the \emph{quasihole degree} $h=q+3-3n$. At the filled flux it equals
$3(N-n)$, three times the electron deficit; a one-hole sector below
means $h=1$, not one missing electron.
Pinning zeros at the
pole realizes the required change of flux. The uniform Fock gap
persists along that pinning deformation and gives local unitary
transport. A history of branch choices then bounds a local
observable by weighted flag values along the history. A first, slowly decaying
weight supplies enough information to close a second estimate with
an integrable weight. Rotation averaging yields
\eqref{eq:intro-zero-bound}.

\paragraph{Local descent of interaction energy.}
Extend the pair Hamiltonian to all particle sectors by
$H_q=\bigoplus_{n=0}^{q+1}H_{n,q}$, with $H_{0,q}=H_{1,q}=0$,
and let $\mathcal N$ act as multiplication by $n$ on the $n$-particle
sector. On this Fock space we construct local particle-loss operators
$J_y$ such that (Theorem~\ref{thm:local-energy-descent})
\begin{equation}\label{eq:intro-descent}
 \int J_y^*J_y\,dy=H_q,\qquad
 \int J_y^*H_qJ_y\,dy\le H_q^2-cH_q.
\end{equation}
Each jump removes between two and a fixed number of particles.
The associated completely positive evolution decreases interaction
energy exponentially, and its expected particle loss is bounded by
the initial energy. The construction first removes a pair and then
partially empties a fixed neighborhood. A rotation average shows
that the remaining pair correlations are paid for by the retained
four-body blocks. This also proves a uniform energy cost for
particle excess above Laughlin filling.

\paragraph{From descent to perturbation theory.}
For a small fixed $\mu>0$, the operator
$G=H_q+\mu(N-\mathcal N)$ controls both interaction energy and
particle-number deficit or excess. If each neutral local term
$A_x$ annihilates $\Omega$, following its expectation through the
loss evolution gives
\[
 \pm\sum_xA_x\le CG.
\]
The proof closes a quadratic inequality for the best relative-bound
constant, using \eqref{eq:intro-zero-bound} at the limiting zero
state (Proposition~\ref{prop:relative-bound}). This implication from
local energy descent and linear electron-deficit
control is useful independently of the particular flux construction.
Finally spectral transport makes the derivative of the pulled-back
perturbed Hamiltonian termwise centered in this sense. On the
$N$-particle sector $G=H_q$, so integration gives a uniform lower
bound above the transported ground state.

Section~\ref{sec:prelim} fixes the algebra and zero-mode conventions.
Sections~\ref{sec:locality}--\ref{sec:charge} establish the local
zero-mode estimate. Section~\ref{sec:energy} constructs the loss
evolution, and Section~\ref{sec:stability} proves the relative bound
and completes Theorem~\ref{thm:main}. Appendix~\ref{app:retained-blocks}
contains the retained-block comparison and its finite arithmetic data.
